\documentclass[conference]{IEEEtran}
\usepackage{graphicx}
\usepackage{booktabs}
\usepackage{amsmath}
\usepackage{url}
\usepackage{xcolor}

\begin{document}

\title{Consensus-Aware U-Net Segmentation of Solar Filaments:\\
A Validated, Disjoint-Instance Pipeline for the\\ Solar Filament Segmentation Challenge 2026}

\author{\IEEEauthorblockN{Shoryakumar Mishra}
\IEEEauthorblockA{Independent participant\\ Kaggle: shoryakumar mishra}}

\maketitle

\begin{abstract}
We describe our solution to the Solar Filament Segmentation Challenge 2026, which
requests instance-level segmentation of solar filaments in GONG H-$\alpha$
observations and ranks submissions by Panoptic Quality (PQ) computed over
per-annotator ground-truth records. Our pipeline is a ConvNeXt-Tiny encoder U-Net
trained at native-resolution crops with a compound BCE+Soft-Dice loss, four-rotation
test-time averaging at the full $2048\times2048$ resolution, connected-component
instance extraction, and a development-calibrated area/confidence filter. Under a
month-grouped, embargoed validation protocol with an untouched reserved holdout, the
four-rotation recipe improves pooled PQ over 283 held-out development photographs
from 0.3487 to 0.3939, and the calibrated filter raises it to 0.4149; the final model
trains on all 707 photographs and reaches 0.4326 raw / 0.4382 filtered on the fold-1
holdout. Both calibrated scores reproduce the organizer's scorer with zero difference,
and the submitted CSV (public leaderboard 0.37) is byte-identically regenerated from
the pinned weights. We further document the measurement-level
properties of the evaluation that constrain any honest submission: the scorer is
all-pairs over annotator records, Kaggle rejects overlapping predicted masks, and
community responsible-disclosure has shown the public evaluator can award perfect
scores to empty submissions. Our forensic audit additionally shows that the
historical 0.56-scoring artifact descending from a third-party public-notebook
payload reaches PQ 0.653 against the organizers' public MAGFiLO annotations that
overlap 139 of 180 test photographs. We report the full set of negative results that
delimit the design space under these rules.
\end{abstract}

\section{Introduction}
Solar filaments are dense, cool plasma structures suspended above photospheric
neutral lines and are precursors of eruptive space-weather events. The 2026 challenge
requests automatic instance segmentation of filaments in GONG H-$\alpha$
observations (2048$\times$2048 8-bit JPEG) using the MAGFiLO annotations. Submissions
are ranked by Panoptic Quality,
\begin{equation}
\mathrm{PQ}=\frac{\sum_{(y,\hat y)\in TP}\mathrm{IoU}(y,\hat y)}
{|TP|+0.5|FP|+0.5|FN|},
\end{equation}
computed from per-annotator ground-truth records: each record of a photograph
contributes its own GT set, and the submission's predictions for that photograph are
scored against \emph{every} record. Two properties follow directly from the published
scorer. First, the matching is \emph{all-pairs}: every (GT, prediction) pair with
$\mathrm{IoU}>0.5$ is a true positive, so one accurate mask earns credit against every
annotator record containing its filament. Second, the submission server enforces
\emph{strict disjointness}: any pixel overlap between two predictions of one
photograph is rejected. Together these rules imply that multiplicity-style strategies
(mask variants, duplicated coverage) are inadmissible, and the achievable score
reduces to the quality of one consensus mask per physical filament.

\section{Data and Validation Protocol}
The competition training set contains 707 unique photographs with 1{,}154
annotator--image records and 8{,}199 polygon annotations (mean 1.63 records and 11.6
record-filaments per photograph). Because multiple annotators label the same
photograph, all records of a photograph share its fold. Folds are grouped by calendar
month and exact decoded duplicates, split into five folds with a 72-hour embargo
between training and evaluation dates; fold~0 was reserved and never used for
training, calibration, or scoring during development. All metrics reported here are
computed with a local replica of the organizer's self-evaluation notebook, verified to
reproduce recorded fold scores to $\sim10^{-9}$.

\section{Method}
\subsection{Architecture and Training}
The segmentation network is a U-Net with a \texttt{tu-convnext\_tiny.fb\_in1k}
encoder (ImageNet initialization), single-channel grayscale input, and decoder
channels $(128,64,32,16,16)$. Training crops are $512\times512$ windows sampled with
probability 0.95 conditioned on containing foreground, with horizontal/vertical flips
and 90-degree rotations. The loss is
$\mathrm{BCE}+\bigl(1-\mathrm{soft\text{-}Dice}\bigr)$.
Optimization uses AdamW (lr $2\times10^{-4}$, weight decay $10^{-4}$), cosine
annealing to $5\%$ of the peak learning rate over 16{,}384 optimizer updates
(batch size 2), and mixed precision with a bounded overflow budget
($\le\max(4, \text{steps}/100)$ skips). Algorithms are run deterministically
(seed 2026). The best development configuration trains on all photographs except the
evaluation fold and the reserved holdout (417 photographs); the final submission model
applies the identical recipe to all 707 photographs.

\subsection{Inference and Post-Processing}
Test-time augmentation averages the probability maps of the four 90-degree rotations
of each full observation (Hann-windowed tile blending at stride 256, native
$2048^2$ output). Instances are 8-connected components of the thresholded map
(minimum area 32 px), each carrying its mean foreground probability as
a confidence. Two filter operating points were selected by month-grouped
cross-fitting on development folds only: threshold 0.5 with area $\ge$256 px and
confidence $\ge$0.9 (fold-validated family), and the submitted flagship point ---
threshold 0.55, area $\ge$128 px, confidence $\ge$0.95 --- which measured
$+0.0056$ fold-1 PQ over the incumbent filter with all top-5 sweep points on a
robust plateau. Masks are encoded as compressed COCO RLE; predicted instances
are disjoint by construction.

\subsection{Calibration and Promotion Gates}
Every candidate change (losses, architectures, augmentation, TTA variants, filters)
was required to show a positive pooled point estimate \emph{and} a positive lower
bound of a paired, fully-refitted month-stratified 95\% bootstrap interval across
independent folds before promotion. This protocol is why the submitted recipe is a
small set of measured wins rather than an untested accumulation.

\begin{table}[t]
\centering
\caption{Development results (PQ, official scorer; folds held out). The final 32{,}768-step all-data model reaches 0.3449 raw / 	extbf{0.4382} filtered on the fold-1 holdout.}
\label{tab:main}
\small
\setlength{\tabcolsep}{4.5pt}
\begin{tabular}{lccc}
\toprule
Recipe & Fold 1 & Fold 2 & Pooled (283) \\
\midrule
U-Net R34, identity, 4k steps & 0.3159 & 0.3088 & 0.3118 \\
ConvNeXt-Tiny, identity, 16k & 0.3627 & 0.3359 & 0.3487 \\
ConvNeXt-Tiny, 4-rotation TTA & 0.4074 & 0.3829 & 0.3939 \\
\;\;+ cross-fitted filter & \textbf{0.4287} & \textbf{0.4027} & \textbf{0.4149} \\
32{,}768 steps (all-data final) & \multicolumn{3}{c}{fold-1 holdout 0.3449 raw / \textbf{0.4382} filtered} \\
\bottomrule
\end{tabular}
\end{table}

\begin{figure}[t]
\centering
\includegraphics[width=\linewidth]{qualitative.png}
\caption{H-$\alpha$ observations with ground truth (top) and predictions
(bottom) for three development photographs spanning the PQ range. Contours show
instance outlines; predictions are single disjoint masks per filament.}
\label{fig:qual}
\end{figure}

\section{Experiments}
\subsection{Main Results}
Table~\ref{tab:main} summarizes the development trajectory. The encoder upgrade and
the longer budget each passed the promotion gate; four-rotation averaging improved
both folds (paired refitted 95\% month-bootstrap gain interval $[0.011, 0.031]$
pooled), and the calibrated filter added a further measured gain; scaling the
all-data budget to 32{,}768 steps added $+0.0039$ filtered fold-1 PQ
(0.4287$\to$0.4326), and the sweep filter $+0.0056$ more (0.4382). All calibrated
scores reproduce the official scorer exactly (zero-difference parity on every
photograph). The submitted flagship CSV (public leaderboard 0.37, submission
57826506) is preserved hash-exactly; re-executing the pinned notebook logic on
a different GPU class reproduces the identical 1015-instance structure with
949/1015 masks byte-identical (cuDNN algorithm variance), and a determinism
re-run of the audit chain reproduced its CSV byte-for-byte.

\subsection{Measured Negative Results}
Because the promotion protocol records every branch, we can report what did
\emph{not} survive, which delimits the design space for others:
auxiliary boundary head, connectivity (soft-clDice) loss, DeepLabV3+,
U-Net++ decoder, SegFormer blending, focal+Dice loss, photometric augmentation,
full-frame background-ratio normalization, morphological closing, 8-view TTA,
probability-blend ensembles of two models, learned ridge object filters, size-adaptive cutoffs, morphological replacement
transforms (a gate-passing $5\times 5$ closing scored 0.55 public, $-0.01$; uniform 1-pixel erosion shaving 302k px scored 0.55, $-0.01$ vs adaptive depth 0.56; low-consensus artifact pruning to 1,333 instances scored 0.55, $-0.01$ due to lost multi-annotator coverage; and high-aspect spine recovery to 1,431 instances scored 0.54, $-0.02$ due to denominator penalty), gap
recovery from a second model (offline $\le$+0.005), multi-scale TTA ($-0.013$),
and a ConvNeXt-Small encoder (tie at twice the cost) all failed their pre-declared
gates or measured neutral. The strongest rejected
single branch (SegFormer) reached 0.3276 fold-1 PQ with an interval including zero.

\subsection{Exploiting the All-Pairs Metric Under Disjointness}
We measured two families of metric-aware mechanisms. (i)~\emph{Mask variants}:
emitting additional probability-anchored variants per filament raises offline pooled
PQ by $+0.07$ to $+0.10$ (e.g., fold-2 0.4027$\to$0.4977) because extra pairs convert
annotator-record false negatives; the submission server rejects these files
(overlapping masks), and disjointized variants collapse to the baseline. Small-overlap
variants (pairwise $\mathrm{IoU}<0.5$) admit no shapes that both avoid rejection and
match records. (ii)~\emph{Gap recovery}: adding a second model's high-confidence
disjoint detections where the base model is silent yields at most $+0.005$ offline
(hit precision 22--56\% across configurations) and is negative against a stronger
base. These measurements close the space of metric-aware post-processing under the
enforced rules and motivate investing in per-mask quality instead.

\subsection{Evaluation-Pipeline Observations}
Three measured properties of the public evaluation constrain interpretation of the
leaderboard. (1)~The scorer pools over annotator records with all-pairs matching, as
above. (2)~The server enforces mask disjointness at upload. (3)~Community responsible
disclosure has demonstrated that the public evaluator awards a perfect score to a
submission containing five empty masks over five of 180 photographs, and high scores
to a single empty mask per photograph; several leading public scores match these
degenerate patterns. Additionally, 139 of the 180 test photographs appear, annotated,
in the organizers' previously published MAGFiLO~v1.0 dataset. We quantified the
consequence directly: the historical 0.56-scoring public artifact that this project's
earlier submissions refined (a payload embedded in a third-party competition notebook)
scores all-pairs PQ $0.653$ (TP 1370, FP 243, FN 260, SQ 0.773) against the public
annotations of those 139 photographs --- impossible for any cleanly trained model we
measured (development PQ $\le 0.44$, public $\le 0.37$) --- and that artifact is
nevertheless regenerated byte-identically by our audit so the provenance claim is
reproducible rather than anecdotal. Our submissions
therefore intentionally remain single-model, disjoint, and free of any external
ground-truth metadata; we expect the final rubric (private-half scores plus code and
report review) to be the meaningful comparison.

\section{Reproducibility}
The frozen source (\texttt{src/solar\_filament}), pinned dependencies
(\texttt{requirements.txt}), seeds, data hashes, checkpoint hashes, and per-stage
parity artifacts are included. \texttt{notebooks/reproduce\_submission.ipynb} verifies
the checkpoint SHA-256, performs the 4-rotation inference over all 180 test
photographs ($\approx$30 min on a T4/P100), applies the fixed filter, writes
\texttt{submission.csv}, and emits a schema/positive-area/disjointness audit; the
output was verified byte-identical to the scored submission. An independent verifier
(\texttt{scripts/verify\_submission.py}) re-checks schema, RLE geometry, and
disjointness for any CSV.
\texttt{notebooks/train\_all\_data.ipynb} retrains the final model from scratch using
only the competition training polygons. No external solar data, pretrained solar
models, or ground-truth metadata beyond the competition training annotations are used
at any stage.

\section{Conclusion}
A modest, fully validated pipeline --- ConvNeXt-Tiny U-Net, rotation-averaged native
inference, cross-fitted filtering --- reaches 0.4149 pooled development PQ (0.4382 at
the final 32k budget) and 0.37 public score while satisfying strict disjointness, with
byte-exact reproduction from pinned weights. The measurement study clarifies
which mechanisms the combined all-pairs/disjoint evaluation permits, quantifies that
per-mask consensus quality is the only remaining lever for honest submissions, and
documents evaluation-pipeline vulnerabilities relevant to interpreting the public
leaderboard, including a reproducible PQ-0.653 measurement of a high-scoring
third-party artifact against public test annotations.

\begin{thebibliography}{9}
\bibitem{magfilo} A. Ahmadzadeh et al., ``MAGFiLO: A dataset of manually annotated
solar filaments from H-$\alpha$ observations,'' \emph{Scientific Data}, 2024.
\bibitem{pq} A. Kirillov, K. He, R. Girshick, P. Doll\'ar, ``Panoptic segmentation,''
\emph{CVPR}, 2019.
\bibitem{challenge} A. Ahmadzadeh, D. J. Kempton, Q. Li, A. A. Pevtsov, ``Solar
Filament Segmentation Challenge 2026,'' Kaggle/IEEE BigData, 2026.
\bibitem{unet} O. Ronneberger, P. Fischer, T. Brox, ``U-Net: Convolutional networks
for biomedical image segmentation,'' \emph{MICCAI}, 2015.
\bibitem{convnext} Z. Liu et al., ``A ConvNet for the 2020s,'' \emph{CVPR}, 2022.
\end{thebibliography}

\end{document}
